\documentclass{article}
\usepackage[utf8]{inputenc}
\usepackage{amsmath}
\usepackage{amsfonts}
\usepackage{booktabs}
\usepackage{hyperref}
\usepackage{geometry}
\usepackage{listings}
\usepackage{xcolor}
\usepackage{fancyhdr}

\geometry{margin=1in}

\lstset{
    basicstyle=\ttfamily\small,
    breaklines=true,
    frame=single,
    backgroundcolor=\color{gray!5},
    showstringspaces=false
}

\fancypagestyle{plain}{%
    \fancyhf{}%
    \chead{\small NATIONAL STUDENT RESEARCH INSTITUTION (NSRI) \\ Open Science $\bullet$ Student Research $\bullet$ 2026}%
    \renewcommand{\headrulewidth}{0.4pt}%
}
\pagestyle{plain}

\title{Majority Illusion in Retrieval-Augmented Generation (RAG)}
\author{
    I. Kamat\textsuperscript{*,1} \and 
    S. Kamsala\textsuperscript{*,1} \and 
    P. Reddi\textsuperscript{*,1} \and 
    K. Sivathenan\textsuperscript{*,1}
}
\date{}

\begin{document}
\maketitle

\begin{center}
    \small
    \textsuperscript{1} Bridgeland High School, Cypress, Texas \\
    \textsuperscript{*} These authors contributed equally to this work. All authors are listed alphabetically.
\end{center}

\begin{abstract}
\textbf{Background:} The Retrieval-Augmented Generation (RAG) model mitigates LLM hallucinations but remains vulnerable to conflicting information in the retrieved context. This study formalizes the ``Majority Illusion'' in RAG, where a model incorrectly adopts a rare or false claim simply because it is over-represented in the retrieved documents.

\textbf{Methods:} We isolate the effect of document frequency from pre-existing model knowledge by evaluating three frontier models (Gemini 3.5 Flash, Claude Haiku 4.5, DeepSeek V4 Flash) on a completely synthetic dataset of 75 fictional entities across varying majority-to-minority document ratios.

\textbf{Results:} We find that all evaluated models exhibit a strong Majority Illusion, consistently adopting the majority claim regardless of its factual validity, though models showed slightly increased caution (flagging conflicts) when handling sensitive financial data.

\textbf{Conclusion:} The reliance on document frequency as a proxy for truth exposes RAG systems to context poisoning attacks. Mitigation requires developing contradiction-aware retrieval mechanisms that prioritize source credibility over frequency.
\end{abstract}

\textbf{Keywords:} Retrieval-Augmented Generation (RAG); Large Language Models; Majority Illusion; In-Context Conflict Resolution; Fictional Entities; System Safety; Chain-of-Thought (CoT) Prompting

\section{Introduction}

The Retrieval-Augmented Generation (RAG) model has become the default design pattern for using Large Language Models (LLMs) in enterprise applications. The use of an external database query in order to create the prompt context allows RAG systems to overcome the drawbacks of parametric memory, such as hallucinations and memory decay. But real-world document databases tend to be dirty, with old records, conflicting reports, repeated extractions, and other noise in them. If LLMs receive conflicting documents from the retrieval process, they have to create context and generate one coherent response out of it.

In the field of network science, the \textbf{majority illusion} is a phenomenon that occurs at the topological level whereby a rare idea or action seems globally prevalent due to its over-representation in the neighborhoods of individual nodes (Lerman \& Ghosh, 2016). We define this principle within the context of RAG systems, whereby in the neighborhood of the retrieved context for the document, a fallacy or outdated information could easily seem to be the general opinion since it is present in many documents. If the LLMs are programmed to solve conflicts using frequency of occurrence in documents, then a heavily repeated and badly curated information can override a true but rare minority claim.

Although there has been some effort towards investigating how LLMs process contradictory context in retrieved settings (e.g., \textit{RamDocs} benchmark (RamDocs, 2025)), there are a number of important issues that have not yet been addressed. The first one is that most tests involve the use of entities/events from the real world. As a consequence, an important confounding variable comes into play in such cases: the prior parametric knowledge of the model (knowledge gained by the model before training). In other words, if the model ``knows'' something, the impact of the contradictory context on the behavior of the model might be the result of prior knowledge and not of document ratios.

In an effort to remedy these limitations, this research employs an artificially constructed evaluation dataset consisting of 75 wholly fabricated entities. Since the entities and facts are artificially constructed, the models will possess no prior parametric information regarding these entities, thus making the reasoning process purely contextual. In order to examine the impact of frequency dominance on model selection, the document ratio is varied across six different conditions, including low, moderate, and high-noise conflicts. This research examines the impact of frequency dominance on three state-of-the-art models and whether CoT prompting serves as a solution to this problem.

\begin{itemize}
    \item \textbf{RQ1 (Majority Dominance):} What is the effect of the majority-to-minority documents' ratio in the prompt context on the likelihood that the model will adhere to the majority claim (\texttt{MAJ}) and on its degree of confidence?
    \item \textbf{RQ2 (Domain Rigidity):} Will domain sensitivity (sensitive financial parameters in banking versus generic corporate facts) affect the caution of the model and increase the number of conflicts that need to be flagged?
    \item \textbf{RQ3 (Position Bias):} Is there an interaction between the spatial position of the minority document (primacy vs. recency) and the majority-to-minority documents' ratio in the prompt context on the output?
\end{itemize}

\section{Methods}

\subsection{Formalization of the Majority Illusion}
The Majority Illusion is defined for a RAG system as follows: let $Q$ be a query of the user and $D = \{d_1, d_2, \dots, d_n\}$ be a collection of documents retrieved by the model in its context window. We split $D$ into two mutually exclusive sets, namely, the majority set $M$ that consists of documents supporting a particular claim $C$, and the minority set $m$ consisting of documents that make a contradictory claim $\neg C$. Denote $|M|$ and $|m|$ as the sizes of these sets, satisfying $|M|+|m|=n$. Then, we define the ratio $r = |M| : |m|$. Majority Illusion happens if the probability of the model outputting the claim $C$, $P(\text{output}=C \mid Q,D)$, tends to 1 while increasing the ratio $r$, irrespective of the validity of $C$. 

In order to observe this phenomenon, we propose an architecture of a prompt which aims to drive the model take the role of synthesizing the documents:
\begin{lstlisting}
System: You are an expert analyst. Answer the user's question 
based STRICTLY on the provided documents. If the documents conflict, 
synthesize the information and state what the documents claim.

Documents:
[Document 1]
...
[Document n]

Question: {Q}
\end{lstlisting}

\subsection{Dataset Generation and Fictional Entities}
In order to control for pre-trained parametric knowledge, we collected a dataset containing 75 artificially generated fictional entities. In order to validate the hypothesis that models have different amounts of caution regarding their predictions depending on sensitivity of the field (Domain Rigidity), the dataset consists of two fields:

\begin{enumerate}
    \item \textbf{Banking (30 entities):} Attributes are strictly numeric sensitive financial parameters, including \texttt{interest\_rate}, \texttt{monthly\_fee}, \texttt{lending\_cap} and \texttt{overdraft\_limit}.
    \item \textbf{General Corporate (45 entities):} Attributes are non-sensitive business facts, including \texttt{founder\_name}, \texttt{headquarters}, \texttt{ceo} and \texttt{employee\_count}.
\end{enumerate}

For each entity, we specify a "majority value" and a conflicting "minority value". After that, we collect a pool of standardized texts using templates (e.g., ``\textit{According to [Source], the CEO of [Entity] is [Value].}''). For each document, we provide a unique mock source (business journals, press releases or financial audits, respectively). The script of dataset generation ensures structural similarity of the context around target attributes in order to exclude style or length bias.

\subsection{Conflicting Document Ratios and Shuffling}
The document context windows are constructed by selecting documents from the pool with specified \textbf{Document Ratios} ($N_{\text{majority}} : N_{\text{minority}}$). We test six ratios:
\begin{itemize}
    \item \textbf{\texttt{4:0} (Control):} 4 majority documents, 0 minority documents. This creates a control for evaluating retrieval accuracy under ideal conditions.
    \item \textbf{\texttt{3:1} (Moderate Conflict):} 3 majority documents, 1 minority document.
    \item \textbf{\texttt{2:2} (Tied Conflict):} 2 majority documents, 2 minority documents.
    \item \textbf{\texttt{4:1} (High Noise Conflict):} 4 majority documents, 1 minority document.
    \item \textbf{\texttt{2:1} (Tied Proportions, Low Count):} 2 majority documents, 1 minority document.
    \item \textbf{\texttt{3:2} (High Noise, Closer Proportions):} 3 majority documents, 2 minority documents.
\end{itemize}

Position bias is counteracted by shuffling the retrieved documents randomly within each trial. This is done using a random seed generated deterministically from the run seed, entity ID, ratio, and trial index:
\begin{equation}
\text{Seed} = f(\text{Run\_Seed}, \text{Entity\_ID}, \text{Ratio}, \text{Trial\_Index})
\end{equation}
This way, we ensure that our experiments are reproducible, yet different trials.

\subsection{Model Evaluation and API Configurations}
The experimental harness asks 3 frontier models using their respective API endpoints:

\begin{itemize}
    \item \textbf{Google Gemini 3.5 Flash:} The model is queried directly using \texttt{google-genai} SDK. In particular, we activate the native capability of Gemini 3.5 Flash using \texttt{thinking\_config} by setting \texttt{thinking\_level = ThinkingLevel.MEDIUM}, since the native capabilities of the model require the user to explicitly activate thinking modes.
    \item \textbf{Claude Haiku 4.5:} Queried using OpenRouter gateway. In order to activate the native reasoning capability of the model, we have to pass \texttt{reasoning: \{"max\_tokens": 2048\}} as part of the \texttt{extra\_body} field of the API request.
    \item \textbf{DeepSeek V4 Flash:} The model is queried using OpenRouter gateway. Just like Claude Haiku, we can activate the reasoning capability by passing \texttt{reasoning: \{"effort": "high"\}} in the \texttt{extra\_body} field.
\end{itemize}

All models are queried with a \texttt{max\_completion} of 2048 tokens to accommodate reasoning traces of large size. Temperature is set to 1.0 to enable non-deterministic sampling, and each condition is sampled 3 times per model to compute the modal response (majority vote across the 3 trials). The harness uses a backoff with jitter for retry in case of transient errors like timeouts or rate limits to avoid data loss due to infrastructure errors.

\subsection{Scoring and Classification Rubric}
The model responses are processed to retrieve the \texttt{answer} and \texttt{probability\_correct} from the JSON block. These responses are then assigned categories based on the following rule-based scoring rubric:
\begin{itemize}
    \item \textbf{\texttt{MAJ}:} The answer is consistent with the majority claim (with minor numeric formatting allowances, such as ``\$12'' being consistent with ``12.00'').
    \item \textbf{\texttt{MIN}:} The answer is consistent with the minority claim.
    \item \textbf{\texttt{COM} (Compromise):} The answer includes or combines elements of both claims (\emph{``The CEO is claimed as Alice by three sources and Bob by one source''}) but does not refuse to choose a value.
    \item \textbf{\texttt{FLAG} (Flagged Conflict):} The model highlights the conflict and refuses to give a choice of value. This category can be identified via regex searches for keywords such as \texttt{conflict}, \texttt{disagree}, \texttt{contradict}, \texttt{inconsistent}, \texttt{cannot be determined}, and \texttt{uncertain}. 
    \item \textbf{\texttt{OTHER}:} The answer is parsed successfully but is not consistent with either claim and does not include any conflict keywords.
    \item \textbf{\texttt{UNSCORED}:} An API error occurred or no valid JSON could be extracted from the completion.
\end{itemize}

\subsection{Statistical Analysis}
To study the majority-follow rate ($p_{\text{MAJ}}$) across different configurations, we compute the proportion of \texttt{MAJ} responses for each model under all ratios. To handle small sample sizes or extreme proportions (0\% or 100\%), we construct 95\% Wilson score confidence intervals, using this formula:
\begin{equation}
\tilde{p} = \frac{p + \frac{z^2}{2n}}{1 + \frac{z^2}{n}} \quad \pm \quad \frac{z}{1 + \frac{z^2}{n}} \sqrt{\frac{p(1-p)}{n} + \frac{z^2}{4n^2}}
\end{equation}
where $p$ is the observed proportion, $n$ is the number of valid trials, and $z = 1.96$.

\section{Results}
Our evaluation across the 75 synthetic entities revealed a pronounced Majority Illusion. In the 4:0 control condition, all models exhibited near-perfect retrieval synthesis, accurately reproducing the context's unified claim. However, as the proportion of conflicting minority documents increased, model behavior diverged sharply. Under the high-noise 4:1 condition, the models demonstrated a strong tendency to adopt the majority claim (\texttt{MAJ}), largely ignoring the single dissenting document. This effect persisted in the 3:1 condition.

Interestingly, domain sensitivity (RQ2) played a measurable role. In the banking domain, which featured strict numerical parameters, the models were marginally more likely to flag conflicts (\texttt{FLAG}) compared to the general corporate domain, suggesting a latent caution heuristic for financial data.

\section{Discussion}
The susceptibility of LLMs to the Majority Illusion introduces critical security vulnerabilities in enterprise RAG systems. Malicious actors could exploit this by executing ``context poisoning'' attacks. By flooding a database with subtly manipulated documents asserting a false narrative, an attacker can ensure that a retrieval system fetches a majority of poisoned documents for specific queries. Because the LLM blindly trusts the majority frequency (e.g., at a 3:1 or 4:0 ratio), it will generate outputs supporting the attacker's narrative. This is particularly dangerous in high-stakes domains such as legal or medical AI assistants, where a user might lack the expertise to independently verify the model's synthesis.

\section{Limitations}
This paper offers empirically strong evidence supporting the phenomenon of the Majority Illusion in RAG systems, but there are limitations to consider. First, concerning \textbf{construct validity} of the current research design, we have employed an artificial data set with 75 entities, which made it possible for us to control the document frequency in isolation from other parameters of the parametric memory. Nevertheless, the real-life situation of document retrieval can imply certain complications in terms of semantic discrepancies between retrieved documents, different document sizes and multi-hop reasoning, and the way in which it affects frequency is unknown to us.

Second, from the point of \textbf{external validity} of the research, we have only examined three frontier commercial models through the use of API. As the models used in the research are proprietary and undergo constant fine-tuning (such as RLHF updates), the conflict-resolution pattern is expected to change over time. Moreover, these findings might not be generalizable to other types of models, such as open-weights (Llama 3) or domain-specific architectures. The future replication of this research in the case of open-weights models is required to carry out an analysis of internal interpretability, including changes in attention weights caused by conflicting tokens.

Regarding \textbf{internal validity}, although we randomly assigned document order in order to counter the effect of primacy and recency bias, our study used a deterministic structure for the prompt. The use of different structures of system prompts such as presenting the model as an adversary to the debater versus the neutral observer may have different intrinsic vulnerability to the illusion. Further research is required to explore this effect.

\section{Conclusion}
This research formalizes and empirically demonstrates the Majority Illusion in Retrieval-Augmented Generation (RAG) systems. Our findings reveal that frontier LLMs are highly susceptible to frequency dominance, often synthesizing a majority narrative even when confronted with conflicting minority evidence. The magnitude of this vulnerability is modulated by the domain of the query (with models showing slightly more caution for sensitive financial data) and the specific ratio of conflicting documents. 

\subsection{Limitations and Future Work}
While this study provides empirical evidence of the Majority Illusion, several limitations must be acknowledged. First, our evaluation was restricted to a specific set of commercial models accessed via APIs. Future work should include open-weights models (e.g., Llama 3) to allow for internal interpretability analysis, specifically examining how attention weights shift during conflicting context scenarios. Second, our synthetic dataset, while designed to isolate the frequency variable, may not fully capture the nuance of real-world, highly complex document conflicts. Future research should investigate mitigation strategies, such as developing contradiction-aware retrieval mechanisms or training models with explicit reinforcement learning protocols to prioritize source credibility over source frequency.

\section*{Abbreviations}
\begin{description}
    \item[API] Application Programming Interface
    \item[CoT] Chain of Thought
    \item[DV] Dependent Variable
    \item[IV] Independent Variable
    \item[LLM] Large Language Model
    \item[RAG] Retrieval-Augmented Generation
    \item[SDK] Software Development Kit
    \item[MAJ] Majority Claim
    \item[MIN] Minority Claim
    \item[COM] Compromise
    \item[FLAG] Flagged Conflict
\end{description}

\section*{Funding}
This research did not receive any specific grant from funding agencies in the public, commercial, or not-for-profit sectors. The authors paid for all API keys and funded their own research.

\section*{Conflicts of Interest}
The authors declare that they have no known competing financial interests or personal relationships that could have appeared to influence the work reported in this paper.

\section*{Ethics Statement}
Ethics approval was not required because the data was fictional.

\section*{Data Availability Statement}
The dataset consists of 75 synthetic data entries designed to simulate RAG environments and is publicly available in our GitHub repository under \texttt{data/entities.json}.

\section*{Authors' Contributions}
[To be completed]

\section*{Acknowledgements}
[To be completed]

\begin{thebibliography}{99}

\bibitem{brown2020}
Brown, T., Mann, B., Ryder, N., Subbiah, M., Kaplan, J. D., Dhariwal, P., ... \& Amodei, D. (2020). Language models are few-shot learners. \textit{Advances in neural information processing systems}, 33, 1877-1901.

\bibitem{carlini2023}
Carlini, N., Ippolito, D., Jagielski, M., Lee, K., Tramer, F., \& Zhang, C. (2023). Extracting training data from large language models. \textit{32nd USENIX Security Symposium (USENIX Security 23)}, 2633-2650.

\bibitem{hasher1977}
Hasher, L., Goldstein, D., \& Toppino, T. (1977). Frequency and the conference of referential validity. \textit{Journal of verbal learning and verbal behavior}, 16(1), 107-112.

\bibitem{ji2023}
Ji, Z., Lee, N., Frieske, R., Yu, T., Su, D., Xu, Y., ... \& Fung, P. (2023). Survey of hallucination in natural language generation. \textit{ACM Computing Surveys}, 55(12), 1-38.

\bibitem{lerman2016}
Lerman, K., Yan, X., \& Wu, X. Z. (2016). The "majority illusion" in social networks. \textit{PloS one}, 11(2), e0147617.

\bibitem{lewis2020}
Lewis, P., Perez, E., Piktus, A., Petroni, F., Karpukhin, V., Goyal, N., ... \& Stoyanov, V. (2020). Retrieval-augmented generation for knowledge-intensive nlp tasks. \textit{Advances in Neural Information Processing Systems}, 33, 9459-9474.

\bibitem{perez2022}
Perez, E., Ringer, S., Lukošiūtė, K., Nguyen, K., Chen, E., Heiner, S., ... \& Kaplan, J. (2022). Discovering language model behaviors with model-written evaluations. \textit{arXiv preprint arXiv:2212.09251}.

\bibitem{touvron2023}
Touvron, H., Lavril, T., Izacard, G., Martinet, X., Lachaux, M. A., Lacroix, T., ... \& Lample, G. (2023). Llama: Open and efficient foundation language models. \textit{arXiv preprint arXiv:2302.13971}.

\bibitem{wallace2019}
Wallace, E., Feng, S., Kandpal, N., Gardner, M., \& Singh, S. (2019). Universal adversarial triggers for attacking and analyzing NLP. \textit{arXiv preprint arXiv:1908.07125}.

\bibitem{wei2023}
Wei, J., Huang, C., Radhakrishnan, A., Wei, L., Pierson, E., \& Steinhardt, J. (2023). Simple synthetic data reduces sycophancy in large language models. \textit{arXiv preprint arXiv:2308.03613}.

\bibitem{wilson1927}
Wilson, E. B. (1927). Probable inference, the law of succession, and statistical inference. \textit{Journal of the American Statistical Association}, 22(158), 209-212.

\end{thebibliography}

\section{Appendix: Prompt Templates}

\subsection{Standard Prompt Template}
\begin{verbatim}
System: You are an expert analyst. Answer the user's question 
based STRICTLY on the provided documents. If the documents conflict, 
synthesize the information and state what the documents claim.

Documents:
{documents}

Question: {question}

Respond with ONLY a JSON object in exactly this format:
{
"answer": "<your answer>",
"probability_correct": <integer from 0 to 100>
}
\end{verbatim}

\subsection{Chain-of-Thought (CoT) Prompt Template}
\begin{verbatim}
System: You are an expert analyst. Answer the user's question 
based STRICTLY on the provided documents. If the documents conflict, 
synthesize the information and state what the documents claim.

Documents:
{documents}

Question: {question}

Before answering, think step-by-step in a "scratchpad" field. 
Analyze the sources, check for conflicts, and determine the most 
accurate synthesis.

Respond with ONLY a JSON object in exactly this format:
{
"scratchpad": "<your step-by-step reasoning>",
"answer": "<your answer>",
"probability_correct": <integer from 0 to 100>
}
\end{verbatim}

\end{document}
